\documentclass[10pt,a4paper]{amsart}
\usepackage{amsmath,amssymb,amsthm,graphicx}
\usepackage{verbatim}
\newenvironment{reference}{\comment}{\endcomment}
%\newenvironment{reference}{}{}
\newenvironment{slogan}{\comment}{\endcomment}
\newenvironment{history}{\comment}{\endcomment}

% For commutative diagrams we use Xy-pic
\usepackage[all]{xy}

% We use 2cell for 2-commutative diagrams.
\xyoption{2cell}
% Theorem environments.
%
\theoremstyle{plain}
\newtheorem{theorem}[subsection]{Theorem}
\newtheorem{proposition}[subsection]{Proposition}
\newtheorem{lemma}[subsection]{Lemma}

\theoremstyle{definition}
\newtheorem{definition}[subsection]{Definition}
\newtheorem{example}[subsection]{Example}
\newtheorem{exercise}[subsection]{Exercise}
\newtheorem{situation}[subsection]{Situation}

\theoremstyle{remark}
\newtheorem{remark}[subsection]{Remark}
\newtheorem{remarks}[subsection]{Remarks}

\numberwithin{equation}{subsection}

% Macros
%
\def\lim{\mathop{\mathrm{lim}}\nolimits}
\def\colim{\mathop{\mathrm{colim}}\nolimits}
\def\Spec{\mathop{\mathrm{Spec}}}
\def\Hom{\mathop{\mathrm{Hom}}\nolimits}
\def\Ext{\mathop{\mathrm{Ext}}\nolimits}
\def\SheafHom{\mathop{\mathcal{H}\!\mathit{om}}\nolimits}
\def\SheafExt{\mathop{\mathcal{E}\!\mathit{xt}}\nolimits}
\def\Sch{\mathit{Sch}}
\def\Mor{\mathop{\mathrm{Mor}}\nolimits}
\def\Ob{\mathop{\mathrm{Ob}}\nolimits}
\def\Sh{\mathop{\mathit{Sh}}\nolimits}
\def\NL{\mathop{N\!L}\nolimits}
\def\CH{\mathop{\mathrm{CH}}\nolimits}
\def\proetale{{pro\text{-}\acute{e}tale}}
\def\etale{{\acute{e}tale}}
\def\QCoh{\mathit{QCoh}}
\def\Ker{\mathop{\mathrm{Ker}}}
\def\Im{\mathop{\mathrm{Im}}}
\def\Coker{\mathop{\mathrm{Coker}}}
\def\Coim{\mathop{\mathrm{Coim}}}

% Boxtimes
%
\DeclareMathSymbol{\boxtimes}{\mathbin}{AMSa}{"02}

%
% Macros for moduli stacks/spaces
%
\def\QCohstack{\mathcal{QC}\!\mathit{oh}}
\def\Cohstack{\mathcal{C}\!\mathit{oh}}
\def\Spacesstack{\mathcal{S}\!\mathit{paces}}
\def\Quotfunctor{\mathrm{Quot}}
\def\Hilbfunctor{\mathrm{Hilb}}
\def\Curvesstack{\mathcal{C}\!\mathit{urves}}
\def\Polarizedstack{\mathcal{P}\!\mathit{olarized}}
\def\Complexesstack{\mathcal{C}\!\mathit{omplexes}}
% \Pic is the operator that assigns to X its picard group, usage \Pic(X)
% \Picardstack_{X/B} denotes the Picard stack of X over B
% \Picardfunctor_{X/B} denotes the Picard functor of X over B
\def\Pic{\mathop{\mathrm{Pic}}\nolimits}
\def\Picardstack{\mathcal{P}\!\mathit{ic}}
\def\Picardfunctor{\mathrm{Pic}}
\def\Deformationcategory{\mathcal{D}\!\mathit{ef}}
\usepackage{cleveref}
\crefname{originallemma}{original lemma}{original lemmas}
\Crefname{originallemma}{Original lemma}{Original lemmas}
\makeatletter
\newlabel{lemma-tautological-same-sheaf}{{00VW}{source}{}{sourcepointer-1}{}}
\newlabel{lemma-tautological-same-sheaf@cref}{{[originallemma][0][]00VW}{[1][0][]source}{}{}{}}
\makeatother
\crefname{originallemma}{original lemma}{original lemmas}
\Crefname{originallemma}{Original lemma}{Original lemmas}
\makeatletter
\newlabel{lemma-compare-sheaf-condition}{{0G1L}{source}{}{sourcepointer-2}{}}
\newlabel{lemma-compare-sheaf-condition@cref}{{[originallemma][0][]0G1L}{[1][0][]source}{}{}{}}
\makeatother
\title{Comparing the sheaves of two sites on the same category}
\author{The Stacks Project Authors}
\date{stacks-project:sites.tex:local-snapshot-20261012; GNU FDL 1.2 or later}
\begin{document}
\maketitle
\noindent\textit{Editorial note.} All delivered mathematical text is the author's own, copied from the pinned Stacks Project chapter file sites-cohomology.tex; no statement or proof step was rewritten, reordered or removed.\par\medskip
\noindent\textit{Editorial note.} The proof is the author's complete argument; every step the author writes out is inside the counted span and no closing sentence is a placeholder.\par\medskip
\noindent\textit{Editorial note.} Every labelled reference occurring inside the counted spans is carried in assembly.external\_references as a non-load-bearing navigation pointer (same-chapter ones at their definition site, cross-chapter ones at the invocation inside the counted span), so no load-bearing mathematics is replaced by a stub. No bibliography citation occurs inside the counted spans.\par\medskip
\noindent\textit{Editorial note.} The checked chapter list for same-name collisions is the set of stacks chapters staged in this repository: chow, cohomology, divisors, etale-cohomology, flat, more-algebra, more-morphisms, morphisms, spaces-more-morphisms, varieties together with this card's chapter sites-cohomology.\par\medskip
\noindent\textit{Editorial note.} The preamble is the chapter's own retained minimal preamble (sources/stacks-sites-cohomology/retained-preamble.tex), a verbatim excerpt of the native preamble: the Xy-pic setup and the native theorem declarations with the mathematical macros.\par\medskip
\noindent\textit{Allowed standard prerequisites (editorial).}\begin{itemize}
\item Tautologically equivalent families give the same sheaf condition. Same chapter, Lemma lemma-tautological-same-sheaf. Supplies the first part of the statement.
\item The sheaf condition is inherited along refinements. Same chapter, Lemma lemma-compare-sheaf-condition. Supplies the passage from the sheaf condition for a refinement to the sheaf condition for the refined covering.
\end{itemize}\medskip
\ifdefined\hypertarget\hypertarget{sourcepointer-1}{}\fi
\noindent\textit{Source reference.} Original lemma 00VW: sources/stacks-sites/pinned-original-source.tex, source lines 1542--1542. Navigation pointer invoked at proof line 1542; the counted span contains the pointer only and the cited result is not reproduced inside it, so the pointer is registered as a non-load-bearing navigation stub in the established form for Stacks material. No bibliography citation is involved.\par\medskip
\ifdefined\hypertarget\hypertarget{sourcepointer-2}{}\fi
\noindent\textit{Source reference.} Original lemma 0G1L: sources/stacks-sites/pinned-original-source.tex, source lines 1556--1556. Navigation pointer invoked at proof line 1556; the counted span contains the pointer only and the cited result is not reproduced inside it, so the pointer is registered as a non-load-bearing navigation stub in the established form for Stacks material. No bibliography citation is involved.\par\medskip
\begin{lemma}
\label{lemma-refine-same-topology}
Let $\mathcal{C}$ be a category. Let $\text{Cov}_i$, $i = 1, 2$
be two sets of families of morphisms with fixed target which
each define the structure of a site on $\mathcal{C}$.
\begin{enumerate}
\item If every $\mathcal{U} \in \text{Cov}_1$ is tautologically
equivalent to some $\mathcal{V} \in \text{Cov}_2$, then
$\Sh(\mathcal{C}, \text{Cov}_2) \subset
\Sh(\mathcal{C}, \text{Cov}_1)$.
If also, every $\mathcal{U} \in \text{Cov}_2$ is tautologically
equivalent to some $\mathcal{V} \in \text{Cov}_1$ then
the category of sheaves are equal.
\item Suppose
that for each $\mathcal{U} \in \text{Cov}_1$ there exists a
$\mathcal{V} \in \text{Cov}_2$ such that $\mathcal{V}$ refines
$\mathcal{U}$. In this case
$\Sh(\mathcal{C}, \text{Cov}_2) \subset
\Sh(\mathcal{C}, \text{Cov}_1)$.
If also for every $\mathcal{U} \in \text{Cov}_2$
there exists a $\mathcal{V} \in \text{Cov}_1$ such that $\mathcal{V}$
refines $\mathcal{U}$, then the categories of sheaves
are equal.
\end{enumerate}
\end{lemma}

\begin{proof}
Part (1) follows directly from Lemma \ref{lemma-tautological-same-sheaf}
and the definitions.

\medskip\noindent
Proof of (2). Let $\mathcal{F}$ be a sheaf of sets for the site
$(\mathcal{C}, \text{Cov}_2)$. Let $\mathcal{U} \in \text{Cov}_1$,
say $\mathcal{U} = \{U_i \to U\}_{i \in I}$. By assumption we may choose a
refinement $\mathcal{V} \in \text{Cov}_2$ of $\mathcal{U}$, say
$\mathcal{V} = \{V_j \to U\}_{j \in J}$ and refinement given
by $\alpha : J \to I$ and $f_j : V_j \to U_{\alpha(j)}$.
Observe that $\mathcal{F}$ satisfies the sheaf condition for
$\mathcal{V}$ and for the coverings $\{V_j \times_U U_i \to U_i\}_{j \in J}$
as these are in $\text{Cov}_2$. Hence $\mathcal{F}$ satisfies
the sheaf condition for $\mathcal{U}$ by
Lemma \ref{lemma-compare-sheaf-condition}.
\end{proof}

\end{document}
